\chapter{نظرية الزمر}\label{ch:b3:groups}

لم تكن الزمر في مجلد السنة الجامعية 2 سوى أدوات للإحصاء: مبرهنة
لاغرانج، والزمر الدائرية، والزمرة المتناظرة وإشارتها. يحوّل هذا
الفصل نظرية الزمر إلى \emph{طريقة}. والمحرّك هو مفهوم زمرة
\emph{تفعل} على مجموعة: فإحصاء المدارات والنقط الصامدة يعطي \omterm{cor:b3:groups:classeq}{معادلة
الأصناف}، ومبرهنة كوشي، ومبرهنات سيلو الثلاث --- وهي المبدأ الأساسي
الذي ينقل الخاصية المحلية إلى الخاصية الشاملة في نظرية الزمر
المنتهية. ثم نتعلّم كيف نركّب الزمر (الجداءات المباشرة ونصف
المباشرة) وكيف نفكّكها (السلاسل التركيبية، والزمر \omterm{def:b3:groups:derived}{القابلة للحل})،
ونبرهن على المبرهنة التي ستُغلق في \cref{ch:b3:galois} مسألة عمرها ثلاثة
قرون حول المعادلات كثيرة الحدود: \omterm{thm:b3:groups:ansimple}{الزمرة المتناوبة} $A_n$ \omterm{def:b3:groups:simple}{بسيطة} من
أجل $n
\geq 5$.

\section{زمر القسمة ومبرهنات التماثل}

في كل ما يلي، $G$ زمرة مكتوبة ضربيًا وعنصرها المحايد $e$. نذكّر بما ورد في مجلد السنة الجامعية 2: الزمر الجزئية،
والصفوف المجاورة $gH$، ومبرهنة لاغرانج ($\abs G = [G:H]\,\abs H$ من أجل $G$
منتهية)، ورتبة عنصر، والزمر الدائرية، والزمرة المتناظرة $S_n$
بتشاكل إشارتها $\varepsilon \colon S_n \to \{\pm1\}$.

\begin{definition}\label{def:b3:groups:normal}
تكون زمرة جزئية $N$ من $G$ \emph{ناظمية}\index{زمرة جزئية ناظمية}
(ونكتب $N \trianglelefteq G$) إذا كان $gNg^{-1} = N$ من أجل كل $g
\in G$ --- وهذا يكافئ أن يتطابق الصفان المجاوران الأيسر والأيمن: $gN =
Ng$ من أجل كل $g$.
\end{definition}

\begin{theorem}[زمرة القسمة]\label{thm:b3:groups:quotient}
لتكن $N \trianglelefteq G$. مجموعة الصفوف المجاورة $G/N$، مزوّدةً بالضرب
$(gN)(hN) = ghN$، زمرةٌ معرَّفة تعريفًا سليمًا، تُسمّى \emph{زمرة
القسمة}\index{زمرة القسمة}، ويكون \emph{الإسقاط القانوني} $\pi \colon G \to G/N$،
$g \mapsto gN$، تشاكلًا شاملًا نواته $N$. وبالعكس، كل نواة تشاكل
زمر \omterm{def:b3:groups:normal}{ناظمية}: فالزمر الجزئية \omterm{def:b3:groups:normal}{الناظمية} هي بالضبط النوى.
\end{theorem}

\begin{proof}
سلامة التعريف هي بيت القصيد. إذا كان $gN = g'N$ و $hN = h'N$،
فلنكتب $g' = gn$ و $h' = hm$ حيث $n, m \in N$. عندئذٍ $g'h' = gnhm =
gh\,(h^{-1}nh)\,m \in ghN$
لأن $h^{-1}nh \in N$ \omterm{def:b3:groups:normal}{بالناظمية}: فجداء الصفين لا يتعلق بالممثّلين
المختارين. أما التجميعية والعنصر المحايد $eN = N$ والمقلوبات $(gN)^{-1} = g^{-1}N$
فموروثة من $G$. ومن الواضح أن $\pi$ تشاكل شامل وأن
$\pi(g) = N \iff g \in N$.

وإذا كان $f \colon G \to H$ تشاكلًا وكان $k \in \ker f$، فإن
$f(gkg^{-1}) = f(g)f(k)f(g)^{-1} = e$: أي أن النوى \omterm{def:b3:groups:normal}{ناظمية}.
\end{proof}

\begin{theorem}[الخاصية الشمولية؛ مبرهنة التماثل الأولى]
\label{thm:b3:groups:firstiso}
ليكن $f \colon G \to H$ تشاكلًا ولتكن $N \trianglelefteq G$ بحيث
$N \subseteq \ker f$. يوجد تشاكل وحيد $\bar f \colon G/N
\to H$ يحقق $f = \bar f \circ \pi$. وعلى وجه الخصوص، بأخذ $N = \ker
f$:\index{مبرهنات التماثل}
\[
G/\ker f \;\xrightarrow{\;\sim\;}\; \operatorname{im} f,
\qquad gN \mapsto f(g).
\]
\end{theorem}

\begin{proof}
الوحدانية: يجب أن يكون $\bar f(gN)$ مساويًا $f(g)$. الوجود: إذا كان
$gN = g'N$ فإن $g^{-1}g' \in N \subseteq \ker f$، ومنه $f(g) = f(g')$، فيكون $\bar
f(gN) = f(g)$
معرَّفًا تعريفًا سليمًا؛ وهو تشاكل لأن $f$ تشاكل. وفي حالة
$N = \ker f$: يكون $\bar f$ متباينًا، لأن $\bar f(gN) = e$ يعني
$g \in \ker f$، أي $gN = N$؛ وصورته هي صورة $f$.
\end{proof}

\begin{theorem}[مبرهنتا التماثل الثانية والثالثة]
\label{thm:b3:groups:secondiso}
لتكن $H \leq G$ ولتكن $N \trianglelefteq G$.
\begin{enumerate}
  \item $HN = \{hn : h \in H,\, n \in N\}$ زمرة جزئية، و $N
        \trianglelefteq HN$، و $H \cap N \trianglelefteq H$، و
        \[
        H/(H \cap N) \;\cong\; HN/N .
        \]
  \item وإذا كان علاوة على ذلك $N \subseteq K \trianglelefteq G$، فإن $K/N
        \trianglelefteq G/N$ و $(G/N)\big/(K/N) \cong G/K$.
\end{enumerate}
\end{theorem}

\begin{proof}
(1) إن $HN$ زمرة جزئية: $(hn)(h'n') = hh'\,(h'^{-1}nh')n' \in HN$
و $(hn)^{-1} = h^{-1}(hn^{-1}h^{-1}) \in HN$، باستعمال \omterm{def:b3:groups:normal}{ناظمية}
$N$. نركّب $H \hookrightarrow HN \xrightarrow{\pi} HN/N$: هذا
التشاكل شامل ($hnN = hN$) ونواته $\{h \in H : h \in
N\} = H \cap N$؛ ثم نطبّق \cref{thm:b3:groups:firstiso}.

(2) الإسقاط $G/N \to G/K$، $gN \mapsto gK$، معرَّف تعريفًا سليمًا
($N \subseteq K$)، وشامل، ونواته $K/N$؛ ونطبّق
\cref{thm:b3:groups:firstiso} مرة أخرى.
\end{proof}

\begin{theorem}[مبرهنة التقابل]\label{thm:b3:groups:correspondence}
لتكن $N \trianglelefteq G$. التطبيق $H \mapsto H/N$ تقابلٌ بين الزمر
الجزئية من $G$ الحاوية على $N$ والزمر الجزئية من $G/N$، وهو يحفظ
الاحتواءات والأدلّة \omterm{def:b3:groups:normal}{والناظمية} (في الاتجاهين).
\end{theorem}

\begin{proof}
تطبيقه العكسي هو $\bar H \mapsto \pi^{-1}(\bar H)$. وكلا التطبيقين يرسل
الزمر الجزئية إلى زمر جزئية، وهما عكسيان أحدهما للآخر: $\pi^{-1}(H/N) =
HN = H$
لأن $N \subseteq H$، و $\pi(\pi^{-1}(\bar H)) = \bar H$
بشمولية $\pi$. ومن الواضح أن الاحتواءات محفوظة؛ و $[G:H] = [G/N : H/N]$ لأن $gH \mapsto (gN)(H/N)$ تقابلٌ
معرَّف تعريفًا سليمًا بين فضاءي الصفوف المجاورة؛ وأخيرًا يكون
$gHg^{-1} = H$ من أجل كل $g$ إذا وفقط إذا كان $(gN)(H/N)(gN)^{-1} = H/N$ من أجل
كل $gN$، وذلك بشمولية $\pi$ أيضًا.
\end{proof}

\begin{example}\label{ex:b3:groups:firstiso}
$\varepsilon \colon S_n \to \{\pm 1\}$ يعطي $S_n/A_n \cong
\{\pm1\}$؛ و $\det \colon GL_n(K) \to K^\times$ يعطي
$GL_n(K)/SL_n(K) \cong K^\times$؛ و $t \mapsto \eu^{2\iu\pi t}$ يعطي
$\R/\Z \cong \mathbb U$، وهي زمرة الدائرة. ومبرهنة التماثل الأولى هي
الوسيلة التي تُحسب بها زمر القسمة \emph{عمليًا}: نبحث عن تطبيق شامل
له النواة المطلوبة.
\end{example}

\begin{method}\label{met:b3:groups:normal}
لإثبات $N \trianglelefteq G$، إليك الطرائق مرتّبة تنازليًا حسب أناقتها:
نُظهر $N$ على أنها نواة تشاكل معرَّف على $G$؛ أو نتحقق من
$gNg^{-1} \subseteq N$ من أجل كل $g$ (وهذا كافٍ: إذ يعطي تطبيقه على
$g^{-1}$ ثم الاقتران الاحتواء المعاكس)؛ أو نتحقق من أن $N$ اتحاد
أصناف اقتران؛ أو نلاحظ أن $[G:N] = 2$ (وعندئذٍ يُفرض $gN = Ng$ ---
\cref{exo:b3:groups:1}).
\end{method}

\section{أفعال الزمر}

\begin{definition}\label{def:b3:groups:action}
\emph{فعل}\index{فعل زمرة} للزمرة $G$ على مجموعة $X$ هو تشاكل
$\varphi \colon G \to \mathfrak{S}(X)$ إلى زمرة تقابلات $X$؛ ونكتب $g \cdot x$ بدل
$\varphi(g)(x)$. وبصيغة مكافئة: تطبيق $G \times X \to X$ يحقق $e \cdot x = x$
و $g
\cdot (h \cdot x) = (gh) \cdot x$. و\emph{مدار}\index{مدار} النقطة
$x$ هو $\mathcal O_x = \{g \cdot x : g \in G\}$، و\emph{مثبِّتها}\index{مثبت} هو
الزمرة الجزئية $G_x = \{g :
g\cdot x = x\}$، و $X^G = \{x : \forall g,\ g \cdot x = x\}$ هي
مجموعة \emph{النقط الصامدة}. ويكون الفعل \emph{متعدّيًا} إذا وُجد
مدار واحد لا غير، و\emph{أمينًا} إذا كان $\varphi$ متباينًا،
و\emph{حرًّا} إذا كانت جميع المثبِّتات تافهة.
\end{definition}

\begin{example}\label{ex:b3:groups:actions}
خمسة أفعال تدير نظرية الزمر المنتهية بأكملها:
\begin{enumerate}
  \item \omterm{def:b3:groups:action}{فعل} $G$ على نفسها عن طريق \emph{الانسحاب الأيسر} $g \cdot x = gx$:
        وهو حر ومتعدٍّ.
  \item \omterm{def:b3:groups:action}{فعل} $G$ على نفسها عن طريق \emph{الاقتران} $g \cdot x = gxg^{-1}$:
        مداراته هي \emph{أصناف الاقتران}\index{صنف اقتران}،
        ومثبِّتاته هي \emph{المُرَكِّزات}
        \index{مركز عنصر} $Z_G(x) = \{g : gx = xg\}$، ونقطه الصامدة
        هي \emph{مركز}\index{مركز زمرة} الزمرة $Z(G)$.
  \item \omterm{def:b3:groups:action}{فعل} $G$ على فضاء الصفوف المجاورة $G/H$ حسب القاعدة $g \cdot xH = gxH$:
        وهو متعدٍّ، ومثبِّت الصف $H$ هو $H$ نفسها. وكل \omterm{def:b3:groups:action}{فعل} متعدٍّ
        من هذا الشكل (\cref{exo:b3:groups:8}).
  \item \omterm{def:b3:groups:action}{فعل} $G$ على مجموعة زمرها الجزئية بالاقتران: مثبِّت $H$ هو
        \emph{الناظِم}\index{ناظم زمرة جزئية} $N_G(H) =
        \{g : gHg^{-1} = H\}$، وهو أكبر
        زمرة جزئية من $G$ تكون $H$ \omterm{def:b3:groups:normal}{ناظمية} فيها.
  \item \omterm{def:b3:groups:action}{فعل} $S_n$ على $\intint{1}{n}$: وهو أمّ الأمثلة جميعًا.
\end{enumerate}
\end{example}

\begin{theorem}[المدار والمثبِّت]\label{thm:b3:groups:orbitstab}
التطبيق $gG_x \mapsto g \cdot x$ تقابلٌ معرَّف تعريفًا سليمًا
$G/G_x \to \mathcal O_x$. وعلى وجه الخصوص، من أجل $G$
منتهية،\index{مبرهنة المدار والمثبت}
\[
\abs{\mathcal O_x} = [G : G_x] \quad\text{يقسم } \abs G,
\]
ولأن المدارات تجزّئ $X$ (فهي أصناف علاقة التكافؤ
$x \sim y \iff y \in \mathcal O_x$)،
\[
\abs X = \sum_{i} \,[G : G_{x_i}]
\qquad (x_i\colon \text{نقطة واحدة في كل مدار}).
\]
\end{theorem}

\begin{proof}
سلامة التعريف والتباين: $gG_x = hG_x \iff h^{-1}g \in G_x \iff
h^{-1}g \cdot x = x \iff g \cdot x = h \cdot x$؛ ويُقرأ التسلسل في
الاتجاهين. أما الشمولية فهي تعريف \omterm{def:b3:groups:action}{المدار} عينه. وتنتج عبارات
الإحصاء من مبرهنة لاغرانج ومن تجزئة $X$ إلى مدارات.
\end{proof}

\begin{corollary}[معادلة الأصناف]\label{cor:b3:groups:classeq}
من أجل زمرة منتهية $G$، وباختيار ممثّل واحد $x_i$ في كل \omterm{ex:b3:groups:actions}{صنف اقتران}
يضم أكثر من عنصر واحد:\index{معادلة الأصناف}
\[
\abs G = \abs{Z(G)} + \sum_i \,[G : Z_G(x_i)],
\qquad\text{كل } [G : Z_G(x_i)] > 1 \text{ يقسم } \abs G.
\]
\end{corollary}

\begin{proof}
نطبّق \cref{thm:b3:groups:orbitstab} على \omterm{def:b3:groups:action}{فعل} الاقتران: فالمدارات الأحادية هي
بالضبط عناصر $Z(G)$.
\end{proof}

\begin{theorem}[النقط الصامدة لزمر $p$]\label{thm:b3:groups:pfixed}
ليكن $p$ عددًا أوليًا. \emph{$p$-زمرة}\index{زمرة p@$p$-زمرة} هي
زمرة منتهية رتبتها قوة للعدد $p$. إذا فعلت $p$-زمرة $G$ على مجموعة
منتهية $X$، فإن
\[
\abs{X^G} \equiv \abs X \pmod p .
\]
ومن نتائج ذلك: كل $p$-زمرة غير تافهة لها مركز غير تافه، وكل زمرة
رتبتها $p^2$ أبيلية.
\end{theorem}

\begin{proof}
عدد عناصر كل \omterm{def:b3:groups:action}{مدار} هو $[G:G_x]$، وهو قوة للعدد $p$؛ وتكون هذه القوة
مساوية $1$ على النقط الصامدة بالضبط، وقابلة للقسمة على $p$ فيما
عدا ذلك. وبالجمع على المدارات نحصل على المطابقة. أما المركز: \omterm{def:b3:groups:action}{ففعل}
اقتران $G$ على نفسها يعطي $X^G = Z(G)$، ومنه $\abs{Z(G)} \equiv
\abs G \equiv 0 \pmod p$،
ووجود $Z(G) \ni e$ يفرض $\abs{Z(G)} \geq
p$. وأما الرتبة $p^2$: فإذا كان
$Z(G) \neq G$ فإن $\abs{Z(G)} = p$ وتكون $G/Z(G)$ دائرية رتبتها
$p$، وهذا يفرض أن تكون $G$ أبيلية (\cref{exo:b3:groups:2}) --- وهو تناقض.
\end{proof}

\begin{theorem}[كوشي]\label{thm:b3:groups:cauchy}
إذا قسم عدد أولي $p$ العدد $\abs G$، فإن $G$ تحتوي على عنصر رتبته
$p$.\index{مبرهنة كوشي}
\end{theorem}

\begin{proof}[برهان (ماكاي)]
لتكن $X = \{(g_1, \dots, g_p) \in G^p : g_1 g_2 \cdots g_p = e\}$.
اختيار $g_1, \dots, g_{p-1}$ اختيارًا حرًّا يحدّد $g_p$: ومنه $\abs X =
\abs G^{p-1}$،
وهو قابل للقسمة على $p$. تفعل الزمرة الدائرية $\Z/p\Z$ على $X$
بالإزاحة الدائرية $(g_1, \dots, g_p) \mapsto (g_2, \dots, g_p,
g_1)$ --- وهذا يحفظ $X$، لأن $g_2 \cdots g_p g_1 =
g_1^{-1}(g_1 \cdots g_p)g_1 = e$.
وحسب \cref{thm:b3:groups:pfixed}،
$\abs{X^{\Z/p\Z}} \equiv \abs X \equiv 0 \pmod p$. والنقط الصامدة هي
المرتّبات الثابتة $(g, \dots, g)$ التي تحقق $g^p = e$؛ والمرتّبة
$(e,
\dots, e)$ إحداها، فعددها $p$ على الأقل، ومنه يوجد على الأقل عنصر
واحد $g \neq e$ يحقق $g^p = e$: ورتبته $p$ بالضبط.
\end{proof}

\begin{theorem}[كايلي]\label{thm:b3:groups:cayley}
كل زمرة رتبتها $n$ تنغمر في $S_n$.\index{مبرهنة كايلي}
\end{theorem}

\begin{proof}
الانسحاب الأيسر $\varphi \colon G \to \mathfrak S(G) \cong S_n$ تشاكل؛ و $\varphi(g) = \mathrm{id}$ يفرض
$g = ge = e$: فهو أمين.
\end{proof}

\begin{method}\label{met:b3:groups:counting}
إحصاء النقط الصامدة هو الافتتاح الشامل لنظرية الزمر المنتهية.
لإثبات \emph{وجود} شيء ما (عنصر مركزي، أو عنصر رتبته $p$، أو \omterm{def:b3:groups:normal}{زمرة
جزئية ناظمية}، أو نقطة صامدة)، نجعل زمرةً مختارةً بعناية تفعل على
مجموعة منتهية مختارة بعناية، ثم نقارن $\abs{X^G}$ مع $\abs X$
بترديد $p$، أو نستفيد من أن أعداد عناصر المدارات تقسم رتبة الزمرة.
وبراهين \cref{thm:b3:groups:pfixed,thm:b3:groups:cauchy} ومبرهنات سيلو الثلاث الآتية خمسُ صياغات
لهذه الفكرة الواحدة.
\end{method}

\section{مبرهنات سيلو}

تقول مبرهنة لاغرانج إن رتبة كل زمرة جزئية تقسم $\abs G$؛ والعكس
خاطئ ($A_4$، ورتبتها $12$، ليس لها زمرة جزئية رتبتها $6$ ---
\cref{exo:b3:groups:1}). تُنقذ مبرهنات سيلو هذا العكس في حالة قوى الأعداد
الأولية، وبندُ الإحصاء فيها هو أمضى أداة عامة نملكها لإنتاج زمر
جزئية \omterm{def:b3:groups:normal}{ناظمية}.

\begin{definition}\label{def:b3:groups:sylow}
نكتب $\abs G = p^a m$ حيث $p \nmid m$. \emph{زمرة سيلو الجزئية
من النمط $p$}\index{زمرة سيلو الجزئية} في $G$ هي زمرة جزئية رتبتها
$p^a$ --- أي $p$-زمرة جزئية أكبر رتبةً ممكنة. ويُرمَز إلى عدد زمر سيلو الجزئية من النمط $p$ في $G$ بالرمز $n_p$.
\end{definition}

\begin{lemma}\label{lem:b3:groups:binom}
إذا كان $\abs G = p^a m$ مع $p \nmid m$، فإن $\dbinom{p^a m}{p^a}
\equiv m \pmod p$.
\end{lemma}

\begin{proof}
في $\mathbb F_p[X]$، يتكرر «حلم المبتدئ» $(1+X)^p = 1 + X^p$ (فالمعاملات
$\binom pk$ من أجل $0<k<p$ قابلة للقسمة على $p$: إذ يقسم $p$ بسط
$\frac{p!}{k!(p-k)!}$ ولا يقسم مقامه) ليعطي $(1+X)^{p^a} = 1 + X^{p^a}$، ومنه
\[
(1+X)^{p^a m} = \bigl(1 + X^{p^a}\bigr)^m
= \sum_{k=0}^{m} \binom mk X^{k p^a} \quad\text{في } \mathbb
F_p[X].
\]
ثم نطابق معامل $X^{p^a}$: فنجد في الطرف الأيسر $\binom{p^a
m}{p^a} \bmod p$، وفي
الطرف الأيمن $\binom m1 = m$.
\end{proof}

\begin{theorem}[سيلو 1: الوجود]\label{thm:b3:groups:sylow1}
من أجل كل عدد أولي $p$، توجد في $G$ زمر سيلو جزئية من النمط $p$.
\end{theorem}

\begin{proof}[برهان (فيلانت)]
لتكن $\Omega$ مجموعة \emph{أجزاء} $G$ التي عدد عناصرها $p^a$؛ تفعل
$G$ على $\Omega$ بالانسحاب الأيسر $g \cdot S = gS$. وحسب
\cref{lem:b3:groups:binom} لدينا $\abs\Omega = \binom{p^a m}{p^a}
\equiv m \not\equiv 0 \pmod p$، ومنه يوجد \omterm{def:b3:groups:action}{مدار} $\mathcal O_S$
عددُ عناصره أولي مع $p$ (إذ لو قسم $p$ عدد عناصر كل \omterm{def:b3:groups:action}{مدار} لقسم
$\abs\Omega$). لتكن $H = G_S$ مثبِّت مثل هذا $S$. وبما أن $[G : H] = \abs{\mathcal O_S}$
أولي مع $p$ وأن $p^a \mid
\abs G = [G:H]\,\abs H$، فإن $p^a \mid \abs H$. وبالعكس، لنثبّت
$s \in S$: فالتطبيق $H \to S$، $h \mapsto hs$، متباين وصورته داخل
$S$ لأن $hS = S$؛ ومنه $\abs H \leq \abs S = p^a$. إذن $\abs H = p^a$.
\end{proof}

\begin{theorem}[سيلو 2: الهيمنة والاقتران]\label{thm:b3:groups:sylow2}
لتكن $P$ زمرة سيلو جزئية من النمط $p$ ولتكن $Q$ أي $p$-زمرة جزئية
من $G$. عندئذٍ $Q \subseteq gPg^{-1}$ من أجل $g \in G$ ما. وعلى وجه الخصوص،
جميع زمر سيلو الجزئية من النمط $p$ مقترنة، و $P \trianglelefteq G \iff
n_p = 1$.
\end{theorem}

\begin{proof}
لتفعل $Q$ على فضاء الصفوف المجاورة $X = G/P$، وعدد عناصره $m
\not\equiv 0 \pmod p$.
وبتطبيق \cref{thm:b3:groups:pfixed} على $p$-زمرة $Q$ نجد
$\abs{X^Q} \equiv m \not\equiv 0 \pmod p$: أي يوجد صف صامد $gP$، أي $QgP = gP$، أي $g^{-1}Qg \subseteq
P$.
وإذا كانت $Q$ نفسها زمرة سيلو جزئية، فتساوي الرتبتين يحوّل
$Q \subseteq gPg^{-1}$ إلى تساوٍ. وأخيرًا يتحقق $P
\trianglelefteq G$ إذا وفقط إذا
تقلّصت مقترناتها $\{gPg^{-1}\}$ --- وهي حسب ما سبق \emph{جميع} زمر
سيلو الجزئية من النمط $p$ --- إلى $\{P\}$.
\end{proof}

\begin{theorem}[سيلو 3: الإحصاء]\label{thm:b3:groups:sylow3}
$n_p \equiv 1 \pmod p$، و $n_p = [G : N_G(P)]$، وهو يقسم $m$.
\end{theorem}

\begin{proof}
لتكن $\mathrm{Syl}_p$ مجموعة زمر سيلو الجزئية من النمط $p$؛ تفعل
$G$ عليها فعلًا متعدّيًا بالاقتران (\cref{thm:b3:groups:sylow2})، ومثبِّت $P$
هو \omterm{ex:b3:groups:actions}{الناظِم} $N_G(P) \supseteq P$: ومنه $n_p = [G :
N_G(P)]$، ويُظهر $m = [G:P] = [G:N_G(P)]\,[N_G(P):P]$ أن
$n_p \mid
m$.

نقصر الآن \omterm{def:b3:groups:action}{الفعل} على $P$ ونحصي النقط الصامدة. إذا كانت $Q \in
\mathrm{Syl}_p$
صامدة تحت \omterm{def:b3:groups:action}{فعل} $P$، فإن $P \subseteq N_G(Q)$؛ وتكون كلٌّ من $P$ و $Q$ زمرة
سيلو جزئية من النمط $p$ في الزمرة $N_G(Q)$، ومنه فهما مقترنتان
فيها (\cref{thm:b3:groups:sylow2} مطبَّقة على $N_G(Q)$)؛ لكن $Q \trianglelefteq N_G(Q)$، ومنه
فإن $Q$ هي مقترنها الوحيد هناك: أي $P = Q$. إذن النقطة الصامدة
الوحيدة هي $P$ نفسها، وتعطي \cref{thm:b3:groups:pfixed} أن $n_p = \abs{\mathrm{Syl}_p}
\equiv \abs{\mathrm{Syl}_p^P} = 1 \pmod p$.
\end{proof}

\begin{method}\label{met:b3:groups:sylowmethod}
لتحليل زمرة رتبتها $n = p^a m$: نسرد قواسم $m$ الموافقة للعدد $1$
بترديد $p$ --- فهذه هي المرشّحة لتكون $n_p$. فإذا كان المرشّح
الوحيد هو $1$، كانت \omterm{def:b3:groups:sylow}{زمرة سيلو الجزئية} من النمط $p$ \omterm{def:b3:groups:normal}{ناظمية}. وإذا
فُرض على $n_p > 1$ أن يكون صغيرًا، فلنفعل بالاقتران على
$\mathrm{Syl}_p$ للحصول على تشاكل $G \to S_{n_p}$ نواته صغيرة.
ولنُحصِ العناصر أيضًا: فزمر سيلو الجزئية المتمايزة ذات الرتبة
\emph{الأولية} $p$ تتقاطع تقاطعًا تافهًا، ومنه فهي تحمل
$n_p(p-1)$ عنصرًا رتبته $p$ بالضبط؛ وكثيرًا ما تفرض الإحصاءات
المتداخلة من أجل أعداد أولية مختلفة تناقضًا (\cref{exo:b3:groups:7}).
\end{method}

\begin{example}\label{ex:b3:groups:pq}
ليكن $\abs G = pq$ حيث $p < q$ عددان أوليان و $p \nmid q - 1$.
عندئذٍ يفرض $n_q \mid p$ و $n_q \equiv 1 \bmod q$ أن $n_q = 1$ (لأن $p <
q$)؛
ويفرض $n_p \mid q$ و $n_p \equiv 1 \bmod p$ أن $n_p = 1$ (لأن $q
\not\equiv 1 \bmod p$).
ولتكن $P$ و $Q$ زمرتَي سيلو الناظميتين: عندئذٍ $P
\cap Q = \{e\}$ (لتباين
الرتبتين)، ومنه $\abs{PQ} = pq$ (\cref{exo:b3:groups:4}) و $G \cong P \times Q \cong \Z/p\Z
\times \Z/q\Z \cong \Z/pq\Z$ حسب
\cref{prop:b3:groups:direct} الآتية. فكل زمرة رتبتها $15$ أو $33$ أو $35$
أو \dots\ دائرية. أما الحالة المستثناة $p \mid q - 1$ فتُنتج زمرة
واحدة إضافية بالضبط، غير أبيلية --- انظر مسألة نهاية الأسبوع
(\cref{pb:b3:groups:1}).
\end{example}

\begin{example}[إحصاء سيلو كاملًا: $S_4$]
\label{ex:b3:groups:s4sylow}
لنطبّق الطريقة على $G = S_4$، حيث $\abs G = 24 = 2^3\cdot3$.
\emph{سيلو $3$:} لدينا $n_3 \mid 8$ و $n_3 \equiv 1 \bmod 3$، ومنه $n_3
\in \{1, 4\}$؛
وبما أن $\langle(123)\rangle$ و $\langle(124)\rangle$ متمايزتان، فإن $n_3 = 4$ --- وهي
الزمر الجزئية الأربع $\langle(abc)\rangle$، واحدة من أجل كل جزء ثلاثي
$\{a, b, c\}$، وهي تستوعب الدورات الثلاثية الثماني. وهي مقترنة
حسب سيلو~2، وتشاكل الاقتران $S_4 \to
S_{\mathrm{Syl}_3} \cong S_4$ تماثلٌ هنا (إذ نواته محتواة
في $N = N_G(\langle(123)\rangle)$ ورتبتها $24/4 =
6$، وكل \omterm{def:b3:groups:normal}{زمرة جزئية ناظمية} من $S_4$
داخل $N$ الشبيهة بالزمرة $S_3$ لا بد أن تكون تافهة: إذ ستتألف من
تبديلات زوجية تُثبّت زمر سيلو الأربع، ولا يفعل ذلك سوى $e$).
\emph{سيلو $2$:} لدينا $n_2 \mid 3$ و $n_2
\equiv 1 \bmod 2$: ومنه $n_2 \in \{1, 3\}$.
والزمرة الجزئية $D =
\langle(1234), (13)\rangle$ رتبتها $8$ (وهي \omterm{ex:b3:groups:semidirectexamples}{زمرة ثنائية السطوح}
$D_4$: تناظرات المربع ذي الرؤوس $1, 2, 3, 4$)، وليست \omterm{def:b3:groups:normal}{ناظمية}
(إذ يولّد $(12)(1234)(12) = (2134)$ زمرة دورات رباعية مغايرة)، ومنه $n_2 = 3$:
فنسخ $D_4$ الثلاث توافق الطرق الثلاث لتزويج $4$ نقط في «مربع».
ولاحظ العبرة من هذا الإحصاء: إن $\abs{S_4} = 24$ يترك مجالًا
لأيٍّ من زمرتَي سيلو ألّا تكون \omterm{def:b3:groups:normal}{ناظمية}، وكلتاهما كذلك \omterm{def:b3:groups:action}{بالفعل} ---
قارن ذلك بالرتبة $12$، حيث يفرض الإحصاء \omterm{def:b3:groups:normal}{ناظمية} إحداهما (الجزء
الرابع من \cref{pb:b3:groups:1}).
\end{example}

\section{الجداءات المباشرة ونصف المباشرة}

\begin{proposition}[التعرّف على جداء مباشر]\label{prop:b3:groups:direct}
لتكن $H, K \trianglelefteq G$ بحيث $H \cap K = \{e\}$ و $HK = G$. عندئذٍ يكون $(h,k) \mapsto hk$
تماثلًا $H \times K \to
G$.\index{جداء مباشر}
\end{proposition}

\begin{proof}
من أجل $h \in H$ و $k \in K$، ينتمي \omterm{def:b3:groups:derived}{المبدِّل} $hkh^{-1}k^{-1}$ إلى
$K$ (بقراءته على الصورة $(hkh^{-1})k^{-1}$، باستعمال \omterm{def:b3:groups:normal}{ناظمية} $K$) وإلى $H$
(بقراءته على الصورة $h(kh^{-1}k^{-1})$): إذن هو $e$، فتتبادل عناصر $H$ مع
عناصر $K$ ويكون التطبيق تشاكلًا. وهو شامل لأن $HK = G$، ومتباين
لأن $hk = e$ يعطي $h = k^{-1} \in
H \cap K = \{e\}$.
\end{proof}

والذي يخفق في أغلب الأحيان هو \omterm{def:b3:groups:normal}{ناظمية} \emph{كلا} العاملين: ففي
$S_3
= \langle (1\,2\,3)\rangle \,\langle(1\,2)\rangle$ يتقاطع العاملان تقاطعًا تافهًا ويولّدان الزمرة، ومع
ذلك $S_3 \not\cong \Z/3\Z \times
\Z/2\Z$. وأما المفهوم الملائم حين يكون عامل واحد ناظميًا
فهو التالي:

\begin{definition}\label{def:b3:groups:semidirect}
لتكن $H$ و $K$ زمرتين وليكن $\varphi \colon K \to
\operatorname{Aut}(H)$ تشاكلًا. \emph{الجداء نصف
المباشر}\index{جداء نصف مباشر} $H \rtimes_\varphi K$ هو المجموعة $H \times K$
مزوّدةً بالقانون
\[
(h, k)\,(h', k') = \bigl(h\,\varphi(k)(h'),\; kk'\bigr).
\]
\end{definition}

\begin{proposition}\label{prop:b3:groups:semidirect}
إن $H \rtimes_\varphi K$ زمرة؛ و $H \times \{e\}$ \omterm{def:b3:groups:normal}{زمرة جزئية ناظمية} تماثل
$H$، و $\{e\} \times K$ زمرة جزئية تماثل $K$؛ وهما تتقاطعان
تقاطعًا تافهًا وتولّدان الزمرة. وبالعكس، إذا كانت $G =
NK$ بحيث
$N \trianglelefteq G$ و $K \leq G$ و $N \cap K = \{e\}$، فإن $G \cong N \rtimes_\varphi K$ حيث $\varphi(k) = (n \mapsto
knk^{-1})$.
\end{proposition}

\begin{proof}
تحقّق مباشر: فالتجميعية ترجع إلى $\varphi(kk') =
\varphi(k)\circ\varphi(k')$ وإلى كون كل $\varphi(k)$
تشاكلًا؛ والعنصر المحايد هو $(e,e)$ و $(h,k)^{-1} =
\bigl(\varphi(k^{-1})(h^{-1}), k^{-1}\bigr)$. والإسقاط $(h,k)
\mapsto k$
تشاكل على $K$ نواته $H \times \{e\}$، ومنه فهذه الأخيرة \omterm{def:b3:groups:normal}{ناظمية}.
أما العكس: فكل $g \in G$ يُكتب \emph{بطريقة وحيدة} على الصورة $nk$
حيث $n \in N$ و $k \in K$ (الوجود: $G =
NK$؛ والوحدانية: $nk = n'k'$
يعطي $n'^{-1}n = k'k^{-1} \in N \cap
K$)، ويُظهر
\[
(nk)(n'k') = n\,(kn'k^{-1})\;kk'
\]
أن $nk \mapsto (n, k)$ ينقل قانون $G$ إلى قانون $N \rtimes_\varphi K$.
\end{proof}

\begin{example}\label{ex:b3:groups:semidirectexamples}
(a) \emph{الزمرة ثنائية السطوح}\index{زمرة ثنائية السطوح} $D_n$
($n \geq
3$) المؤلَّفة من تناظرات المضلع المنتظم ذي $n$ أضلاع وعددها
$2n$: تشكّل الدورانات \omterm{def:b3:groups:normal}{زمرة جزئية ناظمية} دليلها $2$، ويولّد أي
انعكاس متمّمًا لها، ويقلب اقترانُ دورانٍ بانعكاسٍ ذلك الدوران:
$D_n \cong \Z/n\Z \rtimes_\varphi \Z/2\Z$ حيث $\varphi(1) = (x
\mapsto -x)$.
(b) الزمرة التآلفية لمستقيم، $\{x \mapsto ax + b : a \in
K^\times,\, b \in K\} \cong K \rtimes K^\times$: الانسحابات \omterm{def:b3:groups:normal}{ناظمية}،
والتحاكيات متمّم لها.
(c) $S_n \cong A_n \rtimes \Z/2\Z$ (المتمّم: أي مبادلة).
(d) \omterm{pb:b3:groups:1}{زمرة الرباعيات} $Q_8$ \emph{ليست} جداءً نصف مباشر لزمرتين
جزئيتين فعليتين: إذ تحتوي كل زمرة جزئية غير تافهة على $-1$
(\cref{pb:b3:groups:1})، ومنه لا تتقاطع زمرتان جزئيتان فعليتان تقاطعًا
تافهًا.
\end{example}

\section{الزمر القابلة للحل؛ بساطة \texorpdfstring{$A_n$}{An}}

\begin{definition}\label{def:b3:groups:derived}
\emph{المبدِّل}\index{مبدل} للعنصرين $x, y \in G$ هو $[x,y]
= xyx^{-1}y^{-1}$؛
و\emph{الزمرة المشتقة}\index{زمرة مشتقة} $D(G)$ هي الزمرة الجزئية
المولَّدة بجميع المبدِّلات. و\emph{السلسلة المشتقة} هي
$D^0(G) = G$ و $D^{i+1}(G) = D(D^i(G))$، وتكون $G$ \emph{قابلة
للحل}\index{زمرة قابلة للحل} إذا كان $D^n(G) =
\{e\}$ من أجل $n$ ما.
\end{definition}

\begin{proposition}\label{prop:b3:groups:derived}
إن $D(G)$ \omterm{def:b3:groups:normal}{ناظمية} (بل هي صامدة تحت كل تماثل ذاتي)، و $G/D(G)$
أبيلية، ومن أجل $N \trianglelefteq G$: تكون $G/N$ أبيلية $\iff
D(G) \subseteq N$. وعلاوة على
ذلك، تكون $G$ \omterm{def:b3:groups:derived}{قابلة للحل} إذا وفقط إذا وُجدت سلسلة $G = G_0 \trianglerighteq G_1 \trianglerighteq \dots
\trianglerighteq G_n = \{e\}$ بحيث
تكون كل $G_{i+1} \trianglelefteq
G_i$ وكل زمرة قسمة $G_i/G_{i+1}$ أبيلية. والزمر
الجزئية من \omterm{def:b3:groups:derived}{زمرة قابلة للحل} وزمر قسمتها \omterm{def:b3:groups:derived}{قابلة للحل}؛ وبالعكس، إذا
كانت $N$ و $G/N$ قابلتين للحل فكذلك $G$.
\end{proposition}

\begin{proof}
يرسل التماثل الذاتي $\alpha$ العنصر $[x,y]$ إلى $[\alpha x, \alpha y]$: فهو
يبدّل المبدِّلات فيما بينها، ومنه يحفظ الزمرة الجزئية التي تولّدها؛
والاقترانات تماثلات ذاتية، ومن هنا تأتي \omterm{def:b3:groups:normal}{الناظمية}. وفي $G/D(G)$
لدينا $\bar x\bar y \bar x^{-1}\bar y^{-1} = \overline{[x,y]} = \bar e$: \omterm{thm:b3:groups:quotient}{فزمرة القسمة} أبيلية. وإذا كانت $G/N$ أبيلية فإن
كل $[x,y] \in
N$، ومنه $D(G) \subseteq N$؛ وبالعكس، إذا كان $D(G) \subseteq N$ فإن $G/N$،
وهي زمرة قسمة للزمرة الأبيلية $G/D(G)$ حسب مبرهنة التماثل
الثالثة، أبيليةٌ.

فإذا كانت $G$ \omterm{def:b3:groups:derived}{قابلة للحل}، كانت \omterm{def:b3:groups:derived}{السلسلة المشتقة} سلسلةً من هذا
النوع. وبالعكس، إذا أُعطيت سلسلة كهذه، فإن $D^i(G) \subseteq G_i$ بالتراجع: إذ
تعطي أبيلية $G_i/G_{i+1}$ أن $D(G_i) \subseteq G_{i+1}$، ومنه $D^{i+1}(G) = D(D^i G)
\subseteq D(G_i) \subseteq G_{i+1}$؛ وبالتالي
$D^n(G) = \{e\}$.

الوراثة: لدينا $D^i(H) \subseteq D^i(G)$ من أجل $H \leq G$ (بالتراجع)، و $D^i(G/N) = \pi(D^i(G))$
لأن $\pi$ يرسل المبدِّلات على المبدِّلات؛ وهذا يعطي العبارتين
الخاصتين بالزمر الجزئية وزمر القسمة. أما التمديد: فإذا كان
$D^m(G/N) = \{e\}$ فإن $D^m(G) \subseteq N$، ويعطي $D^n(N) = \{e\}$ أن $D^{m+n}(G) = D^n(D^m(G)) \subseteq D^n(N)
= \{e\}$.
\end{proof}

\begin{example}\label{ex:b3:groups:solvableexamples}
الزمر الأبيلية \omterm{def:b3:groups:derived}{قابلة للحل}. و$p$-الزمر \omterm{def:b3:groups:derived}{قابلة للحل}، بالتراجع على
الرتبة: إذ $Z(G) \neq \{e\}$ وتكون $G/Z(G)$ $p$-زمرة أصغر. و $S_3$ و $S_4$
قابلتان للحل: إذ $S_4 \trianglerighteq A_4
\trianglerighteq V \trianglerighteq \{e\}$ حيث $V = \{e,
(1\,2)(3\,4), (1\,3)(2\,4), (1\,4)(2\,3)\}$ هي زمرة كلاين المؤلَّفة من
المبادلات المزدوجة (وهي \omterm{def:b3:groups:normal}{ناظمية} في $S_4$: لأنها اتحاد أصناف
اقتران)، وزمر القسمة أبيلية وهي $\Z/2\Z$ و $\Z/3\Z$ و $V$. وفي
\cref{ch:b3:galois}، ستعني عبارة «المعادلة العامة من الدرجة $n$ \omterm{def:b3:groups:derived}{قابلة للحل}
بالجذور» \emph{حرفيًا} أن «$S_n$ \omterm{def:b3:groups:derived}{زمرة قابلة للحل}». ومن هنا تأتي
أهمية التعريف التالي.
\end{example}

\begin{definition}\label{def:b3:groups:simple}
تكون الزمرة $G \neq \{e\}$ \emph{بسيطة}\index{زمرة بسيطة} إذا لم
تكن لها زمر جزئية \omterm{def:b3:groups:normal}{ناظمية} سوى $\{e\}$ و $G$. والزمرة البسيطة غير
الأبيلية ليست \omterm{def:b3:groups:derived}{قابلة للحل}: إذ إن $D(G) \trianglelefteq G$ ليست $\{e\}$ (وإلا كانت
$G$ أبيلية)، ومنه $D(G) = G$ وتكون \omterm{def:b3:groups:derived}{السلسلة المشتقة} ثابتة. أما
الزمر البسيطة الأبيلية فهي بالضبط الزمر $\Z/p\Z$ حيث $p$ أولي
(فالزمرة الأبيلية بسيطة إذا وفقط إذا لم تكن لها زمرة جزئية فعلية
غير تافهة، أي إذا وفقط إذا كانت دائرية ذات رتبة أولية حسب
لاغرانج).
\end{definition}

\begin{lemma}\label{lem:b3:groups:threecycles}
من أجل $n \geq 3$، تتولّد $A_n$ بالدورات الثلاثية؛ ومن أجل
$n \geq 5$، تكون جميع الدورات الثلاثية مقترنة \emph{داخل $A_n$}.
\end{lemma}

\begin{proof}
كل عنصر من $A_n$ جداءُ عدد زوجي من المبادلات؛ نزاوج بينها ونستعمل
(بالتركيب من اليمين إلى اليسار)
\[
(a\,b)(c\,d) = (a\,c\,b)(a\,c\,d), \qquad
(a\,b)(b\,c) = (a\,b\,c), \qquad
(a\,b)(a\,b) = e
\]
من أجل الأزواج المنفصلة والمتراكبة والمتساوية على الترتيب: فكل
زوج من المبادلات جداءُ دورات ثلاثية.

الاقتران: لدينا $\sigma(a\,b\,c)\sigma^{-1} = (\sigma a\; \sigma b\;
\sigma c)$، ومنه فأي دورتين ثلاثيتين مقترنتان
بواسطة $\sigma \in
S_n$ ما. وإذا كان $\sigma$ فرديًا، فلنستبدل به $\sigma' = \sigma (d\,e)$
حيث $d, e$ نقطتان خارج $\{a, b, c\}$ --- وهما موجودتان لأن
$n \geq 5$؛ عندئذٍ يكون $\sigma'$ زوجيًا و $\sigma'(a\,b\,c)
\sigma'^{-1} = \sigma(a\,b\,c)\sigma^{-1}$، لأن $(d\,e)$
تتبادل مع $(a\,b\,c)$.
\end{proof}

\begin{theorem}[بساطة الزمرة المتناوبة]\label{thm:b3:groups:ansimple}
الزمرة $A_n$ \omterm{def:b3:groups:simple}{بسيطة} من أجل $n \geq 5$.\index{زمرة متناوبة}
\end{theorem}

\begin{proof}
لتكن $N \trianglelefteq A_n$ بحيث $N \neq \{e\}$. حسب \cref{lem:b3:groups:threecycles} يكفي أن نبيّن
أن $N$ تحتوي على دورة ثلاثية \emph{واحدة}: \omterm{def:b3:groups:normal}{فالناظمية} واقتران
الدورات الثلاثية داخل $A_n$ يضعان عندئذٍ جميع الدورات الثلاثية في
$N$، ومنه $N = A_n$.

من أجل $\rho \in S_n$، لتكن $F(\rho) = \{x : \rho(x) \neq x\}$ \emph{حاملها} ولتكن $f(\rho) = \abs{F(\rho)}$.
نختار $\sigma \in N
\setminus \{e\}$ بحيث يكون $f(\sigma)$ \emph{أصغريًا}. ونلاحظ أن كل
تبديلة زوجية غير تافهة تحقق $f \geq 3$، وأن $f(\sigma) =
4$ مستحيل من
أجل $\sigma \in A_n$ إلا إذا كان $\sigma$ مبادلة مزدوجة (فالدورة
الرباعية فردية). وسنبيّن أن $\sigma$ دورة ثلاثية.

\emph{الحالة أ: $\sigma$ جداءُ مبادلات منفصلة}، وليكن $\sigma = (a\,b)(c\,d)\cdots$
حيث $f(\sigma) \geq 4$. نختار $e'
\notin \{a, b, c, d\}$ (وهذا ممكن لأن $n \geq 5$)، ونضع
$\tau =
(c\,d\,e')$ و
\[
\sigma' = \tau\sigma\tau^{-1}\,\sigma^{-1} \in N
\qquad (\tau\sigma\tau^{-1} \in N \text{ بالناظمية}).
\]
وبما أن $\sigma\tau^{-1}\sigma^{-1} = (\sigma c\;\sigma e'\;\sigma
d) = (d\;\sigma e'\;c)$ (باستعمال $\sigma c = d$ و $\sigma d = c$)، نحصل
على $\sigma' = (c\,d\,e')(d\;\sigma e'\;c)$.

إذا كان $\sigma e' = e'$ (وهذا يتحقق خاصةً حين $f(\sigma) =
4$، أي
$\sigma = (a\,b)(c\,d)$): فإن $(d\,e'\,c) =
(c\,d\,e')$ و $\sigma' = (c\,d\,e')^2 = (c\,e'\,d)$، وهي دورة ثلاثية تقع في $N$
وتحقق $f(\sigma') = 3 < 4 \leq f(\sigma)$ --- وهذا يناقض الأصغرية.

وإذا كان $\sigma e' \neq e'$: فإن $\sigma e' \notin \{a, b, c, d,
e'\}$ ($\sigma$ يبادل بين $a,b$ وبين
$c,d$، و $e' \notin \{a,b,c,d\}$ مع تباين $\sigma$)، ومنه يحرّك $\sigma$ النقط
الست $a, b,
c, d, e', \sigma e'$: أي $f(\sigma) \geq 6$. ومن جهة أخرى، فإن $\sigma'$، وهو
جداءُ دورتين ثلاثيتين حاملاهما داخل $\{c, d,
e', \sigma e'\}$، يحقق $f(\sigma') \leq 4$؛
كما أن $\sigma' \neq
e$، لأن $\sigma'(d) = \tau\sigma\tau^{-1}(c) = \tau\sigma(e') =
\sigma e' \neq d$ ($\tau$ يُثبّت $\sigma e' \notin \{c,d,e'\}$). إذن
$\sigma' \in N \setminus\{e\}$ مع $f(\sigma') \leq 4 <
f(\sigma)$: وهذا يناقض الأصغرية.

\emph{الحالة ب: لإحدى دورات $\sigma$ طول $\geq 3$}، وليكن
$\sigma(a) = b$ و $\sigma(b) = c$ حيث $a, b, c$ متمايزة. فإذا كان
$\sigma$ هو هذه الدورة الثلاثية بالضبط، فقد انتهينا. وإلا فإن
$f(\sigma) \geq 5$ (فحالة $f(\sigma) = 4$ مع دورة طولها $\geq3$ هي الدورة
الرباعية الفردية، وهي مستثناة)، ومنه يمكننا أن نختار $d, e' \in
F(\sigma) \setminus \{a, b, c\}$.
نضع $\tau = (c\,d\,e')$ و $\sigma' = \tau\sigma\tau^{-1}\sigma^{-1} \in N$. وكما سبق، فإن $\sigma' = (c\,d\,e')\,(\sigma c\;\sigma e'\;\sigma d)$ لا يحرّك سوى نقط
\[
M = \{c, d, e'\} \cup \{\sigma c, \sigma d, \sigma e'\}
\subseteq F(\sigma)
\]
(فصور النقط المتحرّكة متحرّكة: إذ يعطي $\sigma(x) \ne x$ أن $\sigma(\sigma x) \neq \sigma x$، لتباين
$\sigma$). كذلك $b \notin M$: فالنقط الخمس $a, b, c, d, e'$
متمايزة، ومنه $b
\notin \{c, d, e'\}$؛ ولو كان $b \in \{\sigma c, \sigma d, \sigma
e'\}$ لفُرض $a \in \{c, d, e'\}$ (بتطبيق
$\sigma^{-1}$ واستعمال $\sigma a = b$)، وهذا خاطئ. إذن يُثبّت
$\sigma'$ النقطة $b$ بينما يحرّكها $\sigma$؛ و $F(\sigma') \subseteq F(\sigma)$. وأخيرًا
$\sigma' \neq e$: إذ $\sigma^{-1}(c) = b$ و $\tau^{-1}(b) = b$ و $\sigma(b) = c$ و
$\tau(c) = d$، ومنه $\sigma'(c) = d \neq c$. إذن $\sigma' \in N
\setminus \{e\}$ مع $f(\sigma') \leq f(\sigma) - 1$، وهذا
يناقض الأصغرية.

وبما أن كلتا الحالتين مستحيلة، فإن $\sigma$ دورة ثلاثية.
\end{proof}

\begin{corollary}\label{cor:b3:groups:snnotsolvable}
من أجل $n \geq 5$: لا تكون $A_n$ ولا $S_n$ \omterm{def:b3:groups:derived}{قابلة للحل}، والزمر
الجزئية \omterm{def:b3:groups:normal}{الناظمية} الوحيدة من $S_n$ هي $\{e\}$ و $A_n$ و $S_n$.
\end{corollary}

\begin{proof}
الزمرة $A_n$ \omterm{def:b3:groups:simple}{بسيطة} وغير أبيلية، ومنه فهي غير \omterm{def:b3:groups:derived}{قابلة للحل}
(\cref{def:b3:groups:simple})؛ وكل زمرة تحتوي على زمرة جزئية غير \omterm{def:b3:groups:derived}{قابلة للحل} ليست
\omterm{def:b3:groups:derived}{قابلة للحل} (\cref{prop:b3:groups:derived}). لتكن $N
\trianglelefteq S_n$: عندئذٍ تساوي $N \cap A_n \trianglelefteq A_n$
إما $\{e\}$ وإما $A_n$. فإذا كان $N \cap A_n = A_n$، فإن $A_n \subseteq N$ و $N \in \{A_n, S_n\}$
بالدليل. وإذا كان $N \cap A_n = \{e\}$، فإن قصر الإسقاط $S_n \to S_n/A_n \cong
\Z/2\Z$ على $N$ متباين،
ومنه $\abs N \leq 2$؛ فإذا كان $N = \{e, \sigma\}$، فإن \omterm{def:b3:groups:normal}{الناظمية} تجعل \omterm{ex:b3:groups:actions}{صنف
اقتران} $\sigma$ مساويًا $\{\sigma\}$، أي $\sigma \in Z(S_n)$. لكن
$Z(S_n) = \{e\}$ من أجل $n \geq 3$: إذ لو حرّك $\sigma \ne e$
النقطة $a$ إلى $b \neq a$، فلنختر $c
\notin \{a, b\}$؛ عندئذٍ يرسل $(b\,c)\sigma(b\,c)^{-1}$
النقطة $a$ إلى $c
\ne b$، ومنه فهو يخالف $\sigma$. إذن $N = \{e\}$.
\end{proof}

\begin{theorem}[جوردان--هولدر]\label{thm:b3:groups:jordanholder}
تقبل كل زمرة منتهية $G \neq \{e\}$ \emph{سلسلة
تركيبية}\index{سلسلة تركيبية}
\[
\{e\} = G_0 \trianglelefteq G_1 \trianglelefteq \cdots
\trianglelefteq G_r = G,
\qquad G_{i}/G_{i-1} \text{ بسيطة},
\]
ولا تتعلق المجموعة المتعدّدة \emph{للعوامل التركيبية}
$G_i/G_{i-1}$، إلى حدود التماثل، بالسلسلة المختارة. وتكون الزمرة
المنتهية \omterm{def:b3:groups:derived}{قابلة للحل} إذا وفقط إذا كانت جميع عواملها التركيبية
دائرية ذات رتبة أولية.\index{مبرهنة جوردان-هولدر}
\end{theorem}

\begin{proof}
\emph{الوجود}: بالتراجع على $\abs G$. إذا كانت $G$ \omterm{def:b3:groups:simple}{بسيطة}، أخذنا
$\{e\} \trianglelefteq G$. وإلا اخترنا \omterm{def:b3:groups:normal}{زمرة جزئية ناظمية} فعلية أعظمية $N$ (فعدد
الزمر الجزئية منتهٍ)؛ وتكون $G/N$ \omterm{def:b3:groups:simple}{بسيطة} حسب مبرهنة التقابل (إذ
إن \omterm{def:b3:groups:normal}{زمرة جزئية ناظمية} فعلية غير تافهة من $G/N$ سترتفع إلى \omterm{def:b3:groups:normal}{زمرة
جزئية ناظمية} من $G$ محصورة تمامًا بين $N$ و $G$). ثم نلحق
$N \trianglelefteq G$ \omterm{thm:b3:groups:jordanholder}{بسلسلة تركيبية} للزمرة $N$.

\emph{الوحدانية}: بالتراجع على $\abs G$، وحالة $G$ \omterm{def:b3:groups:simple}{البسيطة}
واضحة. نأخذ سلسلتين تركيبيتين حداهما قبل الأخير $M \trianglelefteq G$ و
$N \trianglelefteq G$ (ومنه $G/M$ و $G/N$ بسيطتان). فإذا كان $M = N$، انتهينا
بالتراجع المطبَّق على $M$. وإلا فإن $MN$، وهي \omterm{def:b3:groups:normal}{ناظمية} في $G$
وتحتوي $M$ احتواءً تامًّا، تساوي $G$ (إذ إن $M$ \omterm{def:b3:groups:normal}{ناظمية} أعظمية:
فأي زمرة \omterm{def:b3:groups:normal}{ناظمية} $M \subsetneq L
\subsetneq G$ سترسَل إلى \omterm{def:b3:groups:normal}{زمرة جزئية ناظمية} فعلية غير
تافهة من $G/M$ \omterm{def:b3:groups:simple}{البسيطة}). وتعطي مبرهنة التماثل الثانية
\[
G/M = MN/M \cong N/(M \cap N),
\qquad
G/N = MN/N \cong M/(M \cap N).
\]
نضع $K = M \cap N$ ($\trianglelefteq G$) ونثبّت \omterm{thm:b3:groups:jordanholder}{سلسلة تركيبية} للزمرة $K$. عندئذٍ
تحمل $M$ سلسلتين تركيبيتين: سلسلتها الأصلية، وسلسلة $K$ متبوعةً
بالحد $K \trianglelefteq
M$ (\omterm{thm:b3:groups:quotient}{فزمرة القسمة} $M/K \cong G/N$ \omterm{def:b3:groups:simple}{بسيطة}). وبالتراجع
(المطبَّق على $M$)، تكون عوامل السلسلة الأصلية للزمرة $M$ هي
$\{\text{عوامل } K\} \cup \{G/N\}$؛ والأمر نفسه من أجل $N$. ومنه فإن لسلسلتَي $G$ العوامل
\[
\{\text{عوامل } K\} \;\cup\; \{\,G/N,\; G/M\,\},
\]
أي المجموعة المتعدّدة نفسها.

القابلية للحل: إذا كانت جميع العوامل من الشكل $\Z/p_i\Z$، كانت
السلسلة سلسلةً ذات زمر قسمة أبيلية، ومنه فإن $G$ \omterm{def:b3:groups:derived}{قابلة للحل}
(\cref{prop:b3:groups:derived}). وبالعكس، فإن كل عامل تركيبي \omterm{def:b3:groups:derived}{لزمرة قابلة للحل} هو
قابل للحل (لأنه زمرة قسمة لزمرة جزئية) وبسيط؛ وكل \omterm{def:b3:groups:simple}{زمرة بسيطة}
\omterm{def:b3:groups:derived}{قابلة للحل} أبيلية (إذ يفرض $D(G) \ne G$ أن $D(G) = \{e\}$)، ومنه
فهي من الشكل $\Z/p\Z$.
\end{proof}

\begin{remark}\label{rem:b3:groups:jhmeaning}
تقول مبرهنة جوردان--هولدر إن كل زمرة منتهية مبنيّة من زمر \omterm{def:b3:groups:simple}{بسيطة}،
مع قائمة أجزاء محدَّدة تحديدًا سليمًا --- أي إنها حسابٌ للزمر تلعب
فيه الزمر \omterm{def:b3:groups:simple}{البسيطة} دور الأعداد الأولية، وتحلّ فيه \emph{كيفيةُ
لصق الأجزاء} (معطيات التمديد، كما في \omterm{def:b3:groups:semidirect}{الجداء نصف المباشر}) محلّ
الضرب المجرّد. أما تصنيف الزمر \omterm{def:b3:groups:simple}{البسيطة} المنتهية --- الزمر الدائرية
$\Z/p\Z$، والزمر المتناوبة $A_{n \geq
5}$، وستّ عشرة عائلة من نمط لي،
و $26$ زمرة شاذة --- فهو أحد معالم رياضيات القرن العشرين؛ وبرهانه،
الممتد على نحو عشرة آلاف صفحة من صفحات الدوريات، يتجاوز هذا
المقرَّر بكثير.
\end{remark}

\begin{omfigure}
\begin{tikzpicture}[xscale=1.55, yscale=1.4]
  \tikzset{sg/.style={draw, rounded corners=2pt, inner sep=3pt,
      fill=omDef!6, font=\small}}
  \node[sg, fill=omThm!10] (G)   at (0, 3)    {$D_4$};
  \node[sg] (K1)  at (-1.6, 2)   {$\langle r^2, s\rangle$};
  \node[sg] (R)   at (0, 2)      {$\langle r\rangle$};
  \node[sg] (K2)  at (1.6, 2)    {$\langle r^2, rs\rangle$};
  \node[sg] (S1)  at (-2.4, 1)   {$\langle s\rangle$};
  \node[sg] (S2)  at (-1.2, 1)   {$\langle r^2 s\rangle$};
  \node[sg, fill=omProp!12] (Z) at (0, 1) {$\langle r^2\rangle$};
  \node[sg] (S3)  at (1.2, 1)    {$\langle rs\rangle$};
  \node[sg] (S4)  at (2.4, 1)    {$\langle r^3 s\rangle$};
  \node[sg] (E)   at (0, 0)      {$\{e\}$};
  \draw (G) -- (K1); \draw (G) -- (R); \draw (G) -- (K2);
  \draw (K1) -- (S1); \draw (K1) -- (S2); \draw (K1) -- (Z);
  \draw (R) -- (Z);
  \draw (K2) -- (Z); \draw (K2) -- (S3); \draw (K2) -- (S4);
  \draw (S1) -- (E); \draw (S2) -- (E); \draw (Z) -- (E);
  \draw (S3) -- (E); \draw (S4) -- (E);
\end{tikzpicture}

{\small الزمر الجزئية العشر \omterm{ex:b3:groups:semidirectexamples}{للزمرة ثنائية السطوح} $D_4 = \langle r,
s \mid r^4 = s^2 = e,\ srs^{-1} = r^{-1}\rangle$. والزمر
الجزئية الثلاث ذات الدليل $2$ (في الصف الأوسط) \omterm{def:b3:groups:normal}{ناظمية}، وكذلك
المركز $\langle r^2\rangle$ (المميَّز بالتظليل)؛ أما زمر الانعكاسات الأربع
فتنقسم إلى صنفَي اقتران من زمرتين. وتعطي السلاسل الممتدة من
الأسفل إلى الأعلى سلاسل تركيبية، مثل $\{e\} \trianglelefteq \langle
r^2\rangle \trianglelefteq \langle r\rangle \trianglelefteq D_4$: وعواملها $\Z/2\Z, \Z/2\Z, \Z/2\Z$
--- وهي دائمًا المجموعة المتعدّدة نفسها، كما تقتضي مبرهنة
جوردان--هولدر.}
\end{omfigure}

\section{تمارين}

\begin{exercise}[$\star$]\label{exo:b3:groups:1}
(a) برهن على أن كل زمرة جزئية دليلها $2$ \omterm{def:b3:groups:normal}{ناظمية}.
(b) برهن على أنه إذا كان $[G : H] = 2$، فإن $x^2 \in H$ من أجل كل
$x \in
G$.
(c) استنتج أن $A_4$ ليس لها زمرة جزئية رتبتها $6$: أي أن عكس
مبرهنة لاغرانج خاطئ. \emph{(أحصِ مربّعات الدورات الثلاثية.)}
\end{exercise}

\begin{exercise}[$\star$]\label{exo:b3:groups:2}
برهن على أنه إذا كانت $G/Z(G)$ دائرية فإن $G$ أبيلية. واستنتج مرة
أخرى أن كل زمرة رتبتها $p^2$ أبيلية، ثم أعطِ، من أجل كل عدد أولي
$p$، مثالًا على زمرة غير أبيلية رتبتها $p^3$. \emph{(فكّر في
المصفوفات المثلثية العليا ذات القطر الواحدي على $\mathbb F_p$.)}
\end{exercise}

\begin{exercise}[$\star$]\label{exo:b3:groups:3}
(a) برهن على أن $\operatorname{Aut}(\Z/n\Z) \cong (\Z/n\Z)^\times$.
(b) برهن على أن التماثلات الذاتية الداخلية $\iota_g \colon x \mapsto
gxg^{-1}$ تشكّل \omterm{def:b3:groups:normal}{زمرة
جزئية ناظمية} $\operatorname{Inn}(G)
\trianglelefteq \operatorname{Aut}(G)$، مع $\operatorname{Inn}(G) \cong G/Z(G)$.
\end{exercise}

\begin{exercise}[$\star\star$]\label{exo:b3:groups:4}
لتكن $H, K$ زمرتين جزئيتين من زمرة منتهية $G$.
(a) برهن على \emph{صيغة الجداء} $\abs{HK}\,\abs{H\cap K} =
\abs H\, \abs K$، بإحصاء ألياف التطبيق
$H \times K
\to HK$، $(h,k) \mapsto hk$.
(b) برهن على أن $HK$ زمرة جزئية إذا وفقط إذا كان $HK = KH$ (وهذا
يتحقق تلقائيًا حين تكون إحداهما \omterm{def:b3:groups:normal}{ناظمية}).
(c) إذا كان $H, K \trianglelefteq G$ و $H \cap K = \{e\}$، فبرهن على أن $hk = kh$ من أجل كل
$h \in H$ و $k \in K$.
\end{exercise}

\begin{exercise}[$\star\star$]\label{exo:b3:groups:5}
(مبرهنة بيرنسايد المساعدة في الإحصاء) لتفعل زمرة منتهية $G$ على
مجموعة منتهية $X$. برهن على أن عدد المدارات هو \emph{المتوسط
الحسابي لعدد النقط الصامدة}:
\[
\#\{\text{مدارات}\} = \frac{1}{\abs G}\sum_{g \in G}
\abs{\operatorname{Fix}(g)},
\qquad \operatorname{Fix}(g) = \{x \in X : g \cdot x = x\},
\]
وذلك بإحصاء المجموعة $\{(g,x) : g\cdot x = x\}$ بطريقتين. تطبيق: تفعل $\Z/p\Z$
($p$ أولي) بالدوران على قلائد مؤلَّفة من $p$ خرزة، مع $a$ لونًا
متاحًا؛ استنتج مبرهنة فيرما الصغرى $a^p \equiv a \pmod p$.
\end{exercise}

\begin{exercise}[$\star$]\label{exo:b3:groups:6}
باستعمال مبرهنات سيلو، برهن على أن كل زمرة رتبتها $15$ دائرية وأن
كل زمرة رتبتها $45$ أبيلية.
\end{exercise}

\begin{exercise}[$\star\star$]\label{exo:b3:groups:7}
برهن على أنه لا توجد \omterm{def:b3:groups:simple}{زمرة بسيطة} رتبتها $30$، ولا \omterm{def:b3:groups:simple}{زمرة بسيطة} رتبتها
$56$. \emph{(من أجل $30$: إذا كان $n_3 \neq 1$ و $n_5 \neq 1$،
فأحصِ العناصر ذات الرتبتين $3$ و $5$. ومن أجل $56$: أحصِ العناصر
ذات الرتبة $7$.)}
\end{exercise}

\begin{exercise}[$\star\star$]\label{exo:b3:groups:8}
(a) لتكن $H \leq G$ ذات الدليل $n$. برهن على أن \omterm{def:b3:groups:action}{فعل} $G$ على $G/H$
يعطي تشاكلًا $G \to S_n$ نواته $\bigcap_{g \in
G} gHg^{-1}$ هي أكبر \omterm{def:b3:groups:normal}{زمرة جزئية ناظمية}
من $G$ محتواة في $H$.
(b) استنتج: إذا كانت $G$ منتهية وكان $p$ \emph{أصغر} قاسم أولي
للعدد $\abs G$، فإن كل زمرة جزئية دليلها $p$ \omterm{def:b3:groups:normal}{ناظمية}.
(c) برهن على أن كل \omterm{def:b3:groups:action}{فعل} متعدٍّ للزمرة $G$ على مجموعة $X$ مماثل
\omterm{def:b3:groups:action}{للفعل} على فضاء صفوف مجاورة: أي يوجد تقابل $X
\to G/G_x$ يتبادل مع
الفعلين.
\end{exercise}

\begin{exercise}[$\star\star$]\label{exo:b3:groups:9}
(a) برهن على أن $D(G)$ هي أصغر \omterm{def:b3:groups:normal}{زمرة جزئية ناظمية} من $G$ ذات زمرة
قسمة أبيلية، وأن كل تشاكل من $G$ إلى زمرة أبيلية يتحلّل بطريقة
وحيدة عبر \emph{الزمرة الأبيلية المرافقة} $G^{\mathrm{ab}} = G/D(G)$.
(b) احسب $D(S_n)$ و $S_n^{\mathrm{ab}}$ من أجل $n \geq 2$، ثم $D(Q_8)$ و
$Q_8^{\mathrm{ab}}$.
\end{exercise}

\begin{exercise}[$\star\star$]\label{exo:b3:groups:10}
لتكن $G$ $p$-زمرة ولتكن $H \subsetneq G$ زمرة جزئية فعلية. برهن
على أن $H \subsetneq N_G(H)$ («النواظم تكبر»)، واستنتج أن كل زمرة جزئية أعظمية
من $p$-زمرة \omterm{def:b3:groups:normal}{ناظمية} ودليلها $p$. \emph{(بالتراجع على $\abs G$،
باستعمال $Z(G) \neq \{e\}$: عالج على حدة الحالتين $Z(G) \subseteq H$ و $Z(G) \not\subseteq H$.)}
\end{exercise}

\begin{exercise}[$\star\star\star$]\label{exo:b3:groups:11}
(بساطة $A_5$، بالممارسة)
(a) برهن على أن أعداد عناصر \omterm{ex:b3:groups:actions}{أصناف الاقتران} في $A_5$ هي $1$ و $15$
و $20$ و $12$ و $12$. وانتبه إلى انشطار صنف الدورات الخماسية في
$S_5$: من أجل دورة خماسية $\sigma$، قارن مُرَكِّزَي $\sigma$ في
$S_5$ وفي $A_5$.
(b) استنتج أن $A_5$ \omterm{def:b3:groups:simple}{بسيطة}: \omterm{def:b3:groups:normal}{فالزمرة الجزئية الناظمية} اتحادُ أصناف
اقتران، وتحتوي على $e$، وعدد عناصرها يقسم $60$.
(c) برهن على أن كل \omterm{def:b3:groups:simple}{زمرة بسيطة} رتبتها $60$ تحقق بالضرورة $n_5 =
6$.
\end{exercise}

\begin{exercise}[$\star\star$]\label{exo:b3:groups:12}
(نواظم زمر سيلو الجزئية ناظمة لنفسها) لتكن $P$ زمرة سيلو جزئية من
النمط $p$ في زمرة منتهية $G$ ولتكن $H =
N_G(P)$.
(a) برهن على أن $P$ هي \omterm{def:b3:groups:sylow}{زمرة سيلو الجزئية} \emph{الوحيدة} من النمط
$p$ في $H$.
(b) استنتج أن $N_G(H) = H$. \emph{(من أجل $g \in N_G(H)$: تكون
$gPg^{-1}$ زمرة سيلو جزئية من النمط $p$ في $H$، ومنه
$gPg^{-1} = P$.)}
(c) استنتج أنه ليس هناك ناظِمُ زمرةِ سيلو محتوًى في \omterm{def:b3:groups:normal}{زمرة جزئية
ناظمية} فعلية من $G$، وأن كل زمرة جزئية أعظمية تحتوي على $N_G(P)$
ناظمةٌ لنفسها.
\end{exercise}

\section{مسألة: الزمر ذات الرتبة 15 على الأكثر}

\begin{problem}[{مسألة نهاية الأسبوع --- تصنيف الزمر
الصغيرة}]\label{pb:b3:groups:1}
الهدف هو تصنيف كامل، ببراهين تامة، للزمر ذات الرتبة $\leq 15$ إلى
حدود التماثل. الرتب $1, 2, 3, 5,
7, 11, 13$ تحسمها مبرهنة لاغرانج (فالزمرة
دائرية)، والرتبتان $4$ و $9$ تحسمهما \cref{thm:b3:groups:pfixed} مع التحليل
الآتي للرتبة $p^2$: فتبقى الرتب $6, 8, 10, 12, 14, 15$.

\medskip
\noindent\textbf{الجزء الأول --- أدوات.}
\begin{enumerate}
  \item برهن على أن كل زمرة يحقق فيها كل عنصر $x^2 =
        e$ أبيليةٌ؛
        واستنتج أن رتبة زمرة منتهية كهذه هي $2^k$ وأنها تماثل
        $(\Z/2\Z)^k$. \emph{(انظر إليها على أنها فضاء متجهي
        على $\mathbb F_2$.)}
  \item برهن على أن كل زمرة رتبتها $p^2$ تماثل $\Z/p^2\Z$
        أو $(\Z/p\Z)^2$. واسرد الزمر الأبيلية ذات الرتبة $8$
        إلى حدود التماثل: $\Z/8\Z$ و $\Z/4\Z \times
        \Z/2\Z$ و $(\Z/2\Z)^3$ ---
        وبرهن على أن القائمة تامة وخالية من التكرار \emph{دون}
        استعمال مبرهنة البنية الواردة في \cref{ch:b3:modules} (ناقش حسب
        الرتبة العظمى لعنصر).
  \item ليكن $\varphi, \varphi' \colon K \to \operatorname{Aut}(H)$ فعلين. برهن على أنه إذا كان $\varphi' = \varphi \circ
        \alpha$ مع
        $\alpha \in \operatorname{Aut}(K)$، فإن $H
        \rtimes_{\varphi} K \cong H \rtimes_{\varphi'} K$.
  \item عيّن $\operatorname{Aut}(\Z/n\Z)$ من أجل $n = 3, 4, 5,
        7$ تعيينًا صريحًا، وبرهن على
        أن $\operatorname{Aut}\bigl((\Z/2\Z)^2
        \bigr) \cong S_3$.
\end{enumerate}

\medskip
\noindent\textbf{الجزء الثاني --- الرتب $2p$ ($6$ و $10$ و $14$)
والرتب $pq$.}
\begin{enumerate}[resume]
  \item ليكن $\abs G = 2p$ حيث $p$ عدد أولي فردي. برهن على أن في
        $G$ \omterm{def:b3:groups:normal}{زمرة جزئية ناظمية} $N = \langle r \rangle$ رتبتها $p$، وعنصرًا
        $s$ رتبته $2$ خارج $N$.
  \item استنتج $G \cong \Z/p\Z \rtimes_\varphi \Z/2\Z$، حيث $\varphi(1) \in \operatorname{Aut}(\Z/p\Z)$ عنصر تقابلي، ثم اخلص إلى:
        $G \cong \Z/2p\Z$ أو $G \cong
        D_p$؛ وتحقق من أن هاتين غير مماثلتين.
        وهذا يحسم الرتب $6$ و $10$ و $14$.
  \item وبصورة أعمّ، ليكن $\abs G = pq$ حيث $p < q$ عددان أوليان.
        برهن على أنه إذا كان $p \nmid q-1$ فإن $G$ دائرية
        (\cref{ex:b3:groups:pq})، وأنه إذا كان $p \mid q - 1$ فتوجد، إضافةً
        إلى $\Z/pq\Z$، زمرة غير أبيلية \emph{واحدة بالضبط}
        $\Z/q\Z \rtimes \Z/p\Z$ إلى حدود التماثل --- استعمل السؤال 3 وكون
        $\operatorname{Aut}(\Z/q\Z) \cong
        (\Z/q\Z)^\times$ دائرية رتبتها $q - 1$، وهي نتيجة نقبلها هنا
        وتُبرهَن في \cref{ch:b3:galois} (دائرية $\mathbb F_q^\times$). ثم اخلص إلى
        نتيجة من أجل الرتبة $15$.
\end{enumerate}

\medskip
\noindent\textbf{الجزء الثالث --- الرتبة $8$.} لتكن $G$ زمرة غير
أبيلية رتبتها $8$.
\begin{enumerate}[resume]
  \item برهن على أن في $G$ عنصرًا $r$ رتبته $4$ (استعمل السؤال 1)
        وأن $N = \langle r\rangle$ \omterm{def:b3:groups:normal}{ناظمية}.
  \item ليكن $s \notin N$. برهن على أن $s^2 \in N$
        (\cref{exo:b3:groups:1}(b))، وأن $srs^{-1} = r^{-1}$ (افحص الصور الممكنة
        للعنصر $r$ بالاقتران، وهي لا بد أن تكون من الرتبة 4، واستبعد
        $srs^{-1} = r$)، وأن $s^2 \in \{e, r^2\}$ \emph{(ماذا يحدث لو كان
        $s^2 = r$ أو $r^3$؟ ولماذا يجب أن يتبادل $s^2$ مع $s$؟)}.
  \item في الحالة $s^2 = e$، برهن على أن $G \cong D_4$.
  \item في الحالة $s^2 = r^2$، برهن على أن جدول الضرب محدَّد
        بالكامل؛ والزمرة الناتجة هي \emph{زمرة
        الرباعيات}\index{زمرة الرباعيات} $Q_8 =
        \{\pm 1, \pm \mathrm i, \pm \mathrm j, \pm \mathrm k\}$، $\mathrm i^2 = \mathrm j^2 = \mathrm k^2 = \mathrm i\,
        \mathrm j\,\mathrm k = -1$
        (ضع $r = \mathrm i$ و $s =
        \mathrm j$). تحقق من وجود $Q_8$،
        مثلًا داخل $GL_2(\C)$ عن طريق
        \[
        \mathrm i \mapsto \begin{pmatrix} \iu & 0\\ 0 & -\iu
        \end{pmatrix},
        \qquad
        \mathrm j \mapsto \begin{pmatrix} 0 & 1\\ -1 & 0
        \end{pmatrix}.
        \]
  \item برهن على أن كل زمرة جزئية غير تافهة من $Q_8$ تحتوي على
        $-1$؛ واستنتج أن كل زمرة جزئية من $Q_8$ \omterm{def:b3:groups:normal}{ناظمية}، وأن
        $D_4 \not\cong Q_8$ (أحصِ العناصر ذات الرتبة $2$)، وأن $Q_8$ ليست
        جداءً نصف مباشر لزمرتين جزئيتين فعليتين.
\end{enumerate}

\medskip
\noindent\textbf{الجزء الرابع --- الرتبة $12$.} ليكن
$\abs G = 12$ و $P_3
\in \mathrm{Syl}_3(G)$ و $P_2 \in \mathrm{Syl}_2(G)$.
\begin{enumerate}[resume]
  \item برهن على أن $n_3 \in \{1, 4\}$ و $n_2 \in \{1, 3\}$، وأن $n_3
        = 4$ يفرض $n_2 = 1$
        \emph{(أحصِ العناصر ذات الرتبة $3$)}.
  \item لنفترض $n_3 = 4$. يعطي \omterm{def:b3:groups:action}{فعل} الاقتران على $\mathrm{Syl}_3$
        تشاكلًا $\rho \colon G \to S_4$. برهن على أن $\ker \rho$، المحتوى في كل
        $N_G(P_3)$ ومنه فرتبته تقسم $3$، تافهٌ \emph{(لماذا لا
        يمكن أن تكون رتبته $3$؟)}؛ وأن الصورة، وهي زمرة جزئية
        رتبتها $12$ من $S_4$، هي بالضرورة $A_4$ \emph{(الزمر
        الجزئية ذات الدليل 2 \omterm{def:b3:groups:normal}{ناظمية} وتحتوي على جميع المربّعات ---
        \cref{exo:b3:groups:1}؛ أحصِ المربّعات في $S_4$)}؛ ثم اخلص إلى
        $G \cong A_4$.
  \item لنفترض $n_3 = 1$، ومنه $G \cong \Z/3\Z \rtimes_\varphi P_2$ مع $\varphi \colon P_2 \to \operatorname{Aut}(\Z/3\Z)
        \cong \Z/2\Z$. اسرد الحالات:
        يعطي $\varphi$ التافه الزمرتين $\Z/12\Z$ و $\Z/6\Z \times \Z/2\Z$؛
        ويعطي $P_2 =
        \Z/4\Z$ مع $\varphi$ شامل \emph{الزمرة ثنائية
        الدائرية} $\mathrm{Dic}_3 = \Z/3\Z \rtimes
        \Z/4\Z$؛ ويعطي $P_2 = (\Z/2\Z)^2$ مع $\varphi$ شامل، إلى
        حدود التكافؤ الوارد في السؤال 3، زمرةً واحدة --- برهن على
        أنها $D_6$، مثلًا بإظهار عنصر رتبته $6$ وعنصر تقابلي شبيه
        بانعكاس.
  \item برهن على أن $\Z/12\Z$ و $\Z/6\Z \times \Z/2\Z$ و $D_6$ و $A_4$ و
        $\mathrm{Dic}_3$ غير مماثلة مثنى مثنى \emph{(أحصِ العناصر
        ذات الرتبة $2$، أو استعمل $n_3$)}. وهذا يحسم الرتبة $12$.
\end{enumerate}

\medskip
\noindent\textbf{الجزء الخامس --- تركيب.}
\begin{enumerate}[resume]
  \item كوِّن جدول التصنيف: من أجل كل رتبة $n \leq
        15$، القائمة
        الكاملة للزمر إلى حدود التماثل، مع الأعداد $1, 1, 1, 2, 1, 2, 1, 5, 2, 2, 1, 5, 1, 2, 1$.
\end{enumerate}

\medskip
\noindent\textbf{الجزء السادس --- إلى ما هو أبعد: الزمر ذات
الرتبة $p^3$ من أجل $p$ فردي.} لتحليل الرتبة $8$ في الجزء الثالث
نظيرٌ بديع من أجل الأعداد الأولية الفردية، مع ظاهرة جديدة تمامًا.
ليكن $p$ عددًا أوليًا فرديًا ولتكن $G$ غير أبيلية رتبتها $p^3$.
\begin{enumerate}[resume]
  \item برهن على أن $\abs{Z(G)} = p$، وأن $G/Z(G) \cong
        (\Z/p\Z)^2$ \emph{(زمرة
        قسمة دائرية على المركز تفرض الأبيلية: \cref{exo:b3:groups:2})}،
        وأن $D(G) =
        Z(G)$ \emph{(من أجل $D(G) \subseteq Z(G)$، استعمل أن $G/Z(G)$
        أبيلية؛ وأما التساوي فلأن $G$ غير أبيلية و $D(G) \neq \{e\}$)}.
        استنتج أن كل مبدِّل $[x, y] = xyx^{-1}y^{-1}$ مركزيٌّ ورتبته تقسم $p$.
  \item (المتطابقة المفتاحية) ليكن $x, y \in G$ وليكن
        $z = [y, x]$، وهو مركزي. برهن بالتراجع على $k$:
        \[
        (xy)^k = x^k y^k z^{k(k-1)/2} .
        \]
        \emph{(انقل كل $y$ عبر كل $x$؛ وكل عبور يكلّف عاملًا
        مركزيًا $z$ واحدًا.)}
  \item استنتج أنه من أجل $p$ فردي يكون التطبيق $\theta\colon x \mapsto
        x^p$ تشاكل
        زمر من $G$ إلى $Z(G)$ \emph{(لماذا يكون $x^p$ مركزيًا؟
        ولماذا تستلزم $z^{p(p-1)/2} = e$ أن يكون $p$ فرديًا؟)}، ثم اخلص إلى
        أن أُسّ $G$ هو $p$ أو $p^2$، وأن الحالتين تتمايزان بحسب
        تفاهة $\theta$.
  \item (الأُسّ $p$) لنفترض أن كل عنصر يحقق $x^p = e$. اختر
        $x, y$ يولّد صنفاهما $G/Z(G)$ وضع $z =
        [y, x]$. برهن على أن
        $z \neq e$، وأن كل عنصر من $G$ يُكتب بطريقة وحيدة على
        الصورة $x^ay^bz^c$ ($0 \leq a, b, c < p$)، وأن الضرب محدَّد بالكامل
        بالعلاقات $x^p = y^p = z^p = e$، و $z$ مركزي، و $[y, x] =
        z$. وتحقق من أن
        \emph{زمرة هايزنبرغ}
        \[
        H_p = \left\{
        \begin{pmatrix} 1 & a & c\\ 0 & 1 & b\\ 0 & 0 & 1
        \end{pmatrix} : a, b, c \in \mathbb F_p \right\}
        \subseteq GL_3(\mathbb F_p)
        \]
        تحقّق هذه العلاقات وأن أُسّها $p$ (احسب $(I + N)^p$ حيث
        $N$ مثلثية عليا تامة، باستعمال $N^3 = 0$ و $p \geq 3$):
        فكل زمرة غير أبيلية رتبتها $p^3$ وأُسّها $p$ تماثل
        $H_p$.
  \item (الأُسّ $p^2$) لنفترض أن لعنصر ما $r \in G$ رتبةً تساوي
        $p^2$، ونضع $N = \langle r\rangle$، وهي \omterm{def:b3:groups:normal}{ناظمية} (لأن دليلها $p$:
        \cref{exo:b3:groups:10}). برهن على أنه يوجد $s \notin N$ يحقق
        $s^p = e$ \emph{(خذ أي $t \notin N$؛ وباستعمال السؤال
        20، صحّحه: $\theta(t) = t^p \in Z(G) \subseteq N$ --- برّر $Z(G) = \langle r^p\rangle$ --- ثم اختر $a$
        بحيث يحقق $s = tr^{a}$ الشرط $s^p = e$؛ وأين استُعملت
        فردية $p$؟)}. وبرهن على أن $srs^{-1} = r^{1+p}$، إلى حدود استبدال $s$
        بقوة له، ثم اخلص إلى أنه توجد زمرة غير أبيلية
        \emph{واحدة} بالضبط رتبتها $p^3$ وأُسّها $p^2$، وهي
        $\Z/p^2\Z \rtimes_\varphi \Z/p\Z$ حيث $\varphi(1)\colon r \mapsto r^{1+p}$ \emph{(استعمل السؤال 3؛ والزمرة
        $\operatorname{Aut}(\Z/p^2\Z)$ دائرية رتبتها $p(p-1)$، وهي نتيجة نقبلها هنا،
        ومنه لها زمرة جزئية وحيدة رتبتها $p$)}.
  \item اخلص إلى الإحصاء: من أجل $p$ فردي توجد $5$ زمر بالضبط
        رتبتها $p^3$ (ثلاث أبيلية واثنتان غير أبيليتين)، تمامًا
        كما في حالة $p = 2$ --- غير أن غير الأبيليتين لم تعودا
        $D_4$ و $Q_8$. وحدّد بدقة موضع انهيار برهان $p$ الفردي
        عند $p = 2$: ففي متطابقة السؤال 19، يساوي $z^{k(k-1)/2}$
        من أجل $k = p = 2$ المقدارَ $z^1 \neq
        e$، فلا يكون التربيع
        تشاكلًا --- وبالفعل، في $Q_8$ عنصر وحيد رتبته $2$ بينما
        وفي $D_4$ ذات الأُسّ $4$ خمسةُ عناصر كهذه.
\end{enumerate}

\medskip
\noindent\textbf{الجزء السابع --- تكملات.}
\begin{enumerate}[resume]
  \item من أجل $p$ فردي، أحصِ العناصر ذات الرتبة $p$ في كلٍّ من
        الزمرتين غير الأبيليتين ذات الرتبة $p^3$: برهن على أن
        في $H_p$ منها $p^3 - 1$ بالضبط، بينما في $M_p =
        \Z/p^2\Z \rtimes \Z/p\Z$ منها
        $p^2 - 1$ بالضبط \emph{(استعمل التشاكل $\theta$ من
        السؤال 20: عيّن صورته، ثم رتبة نواته)}. وتحقق عدديًا من
        أجل $p = 3$: $26$ مقابل $8$. وفسّر لماذا لا يمكن لأي
        حجّة من نوع تشاكل التربيع أن تفصل $D_4$ عن $Q_8$، وأيُّ
        إحصاء يفصل بينهما.
  \item نقول عن عدد صحيح $n \geq 1$ إنه \emph{دائري} إذا كانت كل
        زمرة رتبتها $n$ دائرية. برهن على أنه إذا كان $p^2 \mid n$
        من أجل عدد أولي $p$ ما، أو إذا كان للعدد $n$ قاسمان أوليان
        $p < q$ بحيث $p \mid q - 1$، فإن $n$ ليس دائريًا \emph{(في
        كل حالة أعطِ زمرة غير دائرية رتبتها $n$، مستعملًا الجزء
        الثاني في الحالة الثانية)}. استنتج أن دائرية $n$ تفرض
        $\gcd(n, \varphi(n)) = 1$، حيث $\varphi$ دالة أويلر، وتحقق من ذلك في جدول
        السؤال 17: فمن بين $n \leq 15$، الرتب التي تحمل زمرة
        واحدة هي بالضبط $n \in \{1, 2, 3, 5, 7, 11, 13, 15\}$، وهي بعينها الرتب التي تحقق
        $\gcd(n, \varphi(n)) = 1$.
\end{enumerate}
\end{problem}
